%% ============================================================
%% sections/a7-explainability.tex
%% H. R. 9510 Bill v3.0 - APPENDIX C. THE VVUQ EXPLAINABILITY STANDARD
%% (NON-OPERATIVE). Establishes the VVUQ-01 through VVUQ-04 developments as a new
%% standard for explainability of the U.S. legislation-creation process.
%% Common visuals: Figure 9 (prompt evolution across the works), Figure 10
%% (internal bill evolution), Table 7 (the explainability standard).
%% ============================================================
\clearpage
\phantomsection
\addcontentsline{toc}{section}{Appendix C. The VVUQ Explainability Standard (Non-Operative)}
\usccrosshead{APPENDIX C. THE VVUQ EXPLAINABILITY STANDARD (NON-OPERATIVE)}

\uscprov{1}{This appendix is non-operative. It establishes the VVUQ-01 through
VVUQ-04 developments as a reusable standard for explaining how a piece of United
States legislation was created. The claim is narrow and demonstrable: every step
that produced this bill, from the first method paper to the precise amendment,
recorded the exact prompt that drove it, the full output it produced, a VVUQ gate
decision that held the output to its specification, and a commit trail that pushed
each file in real time. A reader can therefore reconstruct how the law was built,
which is the property an explainable process must have \cite{nist-ai-explainability,
nist-ai-600}.}

\figslot{Figure 9.}{Prompt Evolution Across the Works}{The chain of recorded
prompts from VVUQ-02 and VVUQ-03 (instruct, draft, full) through VVUQ-04 (instruct,
update, draft, full) to VVUQ-05 (figures, next-steps, this draft), showing how the
instruction to the agent sharpened as the work moved from method to statute.}

\uscprov{1}{The prompt record is the first half of the standard. Because each
stage keeps its prompt and its output side by side, the reasoning that turned
evidence into statutory language is inspectable rather than hidden: a reviewer can
read prompt-a-update-bill and see the instruction that retargeted the bill onto
the Federal Food, Drug, and Cosmetic Act, then read the output that resulted
\cite{vvuq05-figures-bill}.}

\figslot{Figure 10.}{Internal Bill Evolution}{The VVUQ-04 internal stages
instruct-bill, template-bill, update-bill, draft-bill, full-bill, and final-bill,
extended by the VVUQ-05 update-bill (figures-bill and next-steps) into this v3.0
draft-bill, each stage a labeled box feeding the next.}

\uscprov{1}{The stage record is the second half. The bill did not appear in one
pass: a legal crosswalk was assembled, the real statute was reproduced, the
amendment was scaffolded with bracketed instructions, the instructions were
executed, and the result was finalized, with this visual draft a further stage.
Each stage is a directory a reviewer can open, so the provenance of every clause
is traceable to the files it was synthesized from \cite{kawchak2026vvuq04bill}.}

\figslot{Table 7.}{The VVUQ Explainability Standard}{Each explainability dimension
(explanation, meaningful, accuracy, knowledge limits), what makes the
legislation-creation step satisfy it, and the artifact in the repository that
records it.}

\uscprov{1}{Mapped to the four principles of explainable artificial intelligence,
the standard reads as follows: explanation is met because every step has a
recorded prompt and output; meaningful is met because each stage carries a README
that explains the step to a human; accuracy is met because the VVUQ gate holds
each output to its specification and the results reproduce from a fixed seed; and
the knowledge-limits principle is met because the gate escalates to a human and the
disclaimers mark where the system must hand back. This is the explainability the
legislation-creation process should carry, demonstrated end to end
\cite{nist-ai-explainability, kawchak2026vvuq01}.}

\begin{draftbox}{Appendix C - the explainability standard and its two evolution figures}
\dstep{Figure 9 (prompt evolution): render a left-to-right box-and-arrow chain in the asciifig environment from the recorded prompt files, in this order: papers/VVUQ-02/instructions/prompt-instruct.md, papers/VVUQ-02/draft-paper/prompt-draft-paper.md, papers/VVUQ-02/full-paper/prompt-full-paper.md; papers/VVUQ-03/instruct-paper/prompt-instruct-paper.md, papers/VVUQ-03/draft-paper/prompt-draft-paper.md, papers/VVUQ-03/full-paper/prompt-full-paper.md; papers/VVUQ-04/instruct-bill/prompt-instruct-bill.md, papers/VVUQ-04/update-bill/prompt-update-bill.md and prompt-a-update-bill.md, papers/VVUQ-04/draft-bill/prompt-draft-bill.md, papers/VVUQ-04/full-bill/prompt-full-bill.md; and papers/VVUQ-05/update-bill/figures-bill/prompt-figures-bill.md, papers/VVUQ-05/update-bill/next-steps/prompt-1-next-steps.md and prompt-2-next-steps.md, papers/VVUQ-05/draft-bill/prompt-draft-bill.md. Label each box with the stage and a few words of how the instruction sharpened; keep exact spacing on a white background.}
\dstep{Figure 10 (internal bill evolution): render a box-and-arrow chain in the asciifig environment over the VVUQ-04 stage directories papers/VVUQ-04/instruct-bill, template-bill, update-bill, draft-bill, full-bill, and final-bill, extended by papers/VVUQ-05/update-bill (figures-bill and next-steps) to papers/VVUQ-05/draft-bill (this v3.0 draft). Reuse the lineage motif from papers/VVUQ-05/update-bill/figures-bill/05VVUQ04.md ("Purpose and place in the lineage") and papers/VVUQ-04/final-bill/README.md ("Lineage"); mark the v3.0 endpoint with a "<== THIS DRAFT" pointer.}
\dstep{Table 7 (the explainability standard): build a left-aligned, ragged-right table at the body measure (each width prepended with \rabs) with columns Dimension, What makes the legislation-creation step satisfy it, and Recording artifact. Use the four NISTIR 8312 principles (explanation, meaningful, accuracy, knowledge limits) for the rows; resolve to references.bib keys nist-ai-explainability and nist-ai-600. Map each row to the concrete artifacts: the prompt and output Markdown pairs at each stage; the per-stage README; the VVUQ gate decision surface and the seed-20260525 reproducibility; and the ESCALATE rule and the disclaimers.}
\dstep{Framing source: ground the standard in papers/VVUQ-05/update-bill/figures-bill/README.md (Section 4, how the five catalogs correlate, and Section 5, how a downstream AI should use the corpus) and papers/README.md ("The processing feat accomplished by AI"), which already describe the self-similar, fully recorded pipeline. Keep the claim narrow and demonstrable: this is an explainability standard for the legislation-creation process, not a claim about the merits of the bill.}
\dstep{Formatting: single hyphens only; the section symbol where a codified reference appears; even RaggedRight spacing with no overflow; keep both figures and the table on a white background; size each asciifig to its content so no large empty white band opens, and avoid leaving a single line stranded onto the next page.}
\end{draftbox}
